% DRAFT (LLM-drafted from the protocol documents; the author reviews and confirms).
\section{Pre-registration chain and request manifests}
\label{app:protocol}

\paragraph{Headline run.} The protocol and its addendum were committed before any code that could score the test
set; the selection of \orb's configuration was a separate commit containing only out-of-fold scores on the
selection split; the single test run was the only commit to touch the test-time prediction file, and its metadata
records that it was not re-run. Commit identifiers are withheld for review; the protocol documents in the
supplementary archive record the order and dates of these commits, and the de-anonymised version gives the
identifiers.

\paragraph{Thirteen-task study.} Its protocol (\S\ref{sec:res-13task}) and nine amendments were pre-registered in a
private repository, each before the runs it governs. The released copy is a de-identified transcription made after
the runs: host names, internal paths, file names of operational scripts and names of unrelated services were removed
or replaced by a neutral description; no rule, number, threshold or timestamp was changed.

\paragraph{Closed-API runs.} The 24-shot swap protocol, its addendum and the analysis script were committed before
the first scored call; the zero-shot protocol and its addendum likewise. Every request body's hash is recorded in a
manifest, every call's outcome in a ledger, and the served model version with every response; \fsJevFsCalls{} 24-shot calls succeeded on the first attempt, and the
zero-shot run needed \fsJevZsFailedAttempts{} retried attempt with no failed item.

\paragraph{Open-model runs.} Each open-model run was pre-registered with amendments before the runs they govern
(new engine, new cells). The serving-engine parity gate compares, item by item on the full Banking77 test set, the
engine's probabilities with the release code's (argmax agreement and maximum absolute difference).

\paragraph{Request formats.} Every system received the same JSON body; the open models ignore the
\texttt{model} field. The 24-shot bodies follow the reproduction's format: its instruction and its one-line
definition of each of the 77 intents, which we read from the reproduction at run time and do not redistribute
because it carries no licence, and the retrieved examples in rank order, trimmed from the end if the serialised body
exceeds 30,000 bytes (the reproduction's rule). Schematically: % numbers-ok: the reproduction's byte cap
\begin{verbatim}
{"model": "<pinned version>",
 "state": {"customer_message": "<test utterance>",
           "labeled_examples": [{"message": "<retrieved text>", "intent": "<label>"},
                                ... 24 in rank order ...]},
 "questions": {"intent": {"type": "choice",
                          "instructions": "<the reproduction's instruction>",
                          "criteria": {"<label>": "<the reproduction's definition>",
                                       ... 77 labels ...}}}}
\end{verbatim}
The Bitext bodies have the same structure with our own instruction (``Which ONE intent best describes the customer
message? The labeled examples show messages with their intents. Every message belongs to exactly one of the listed
intents.'') and the one-line descriptions of \S\ref{sec:systems} as criteria. The zero-shot bodies omit
\texttt{labeled\_examples} and use either the label names or the descriptions as criteria. The thirteen-task study
used the bodies built by each dataset's builder or by the workflow harness, recorded by hash in its manifests. The
released manifests list, for every body, its SHA-256, its size, the number of examples kept and the indices of the
retrieved examples, so that every body can be rebuilt and checked.

\begin{table}[h]
\centering\small
\caption{Versions of every system and component, as recorded in the run metadata.}
\label{tab:versions}
\begin{tabular}{@{}lll@{}}
\toprule
System or component & Identifier & Revision \\
\midrule
\jev{} & vendor API, model field & \texttt{jev-1.13.0} (every response) \\ % numbers-ok: version string
\clef{} & open weights, release code & \texttt{2f3de3d} \\ % numbers-ok: revision hash
\flash{} & open weights, release code & \texttt{17f0b0a} \\ % numbers-ok: revision hash
\pplx{} & open weights, v1.1 & \texttt{3b45dea} \\ % numbers-ok: revision hash
Serving engine & SGLang container image & digest \texttt{868b0bd8\ldots} \\ % numbers-ok: image digest
\orb{} reader & Qwen3-4B & \texttt{1cfa9a7} \\ % numbers-ok: model name and revision hash
\orb{} reader adapter & ours (released) & withheld for review \\
\orb{} kNN member & bge-large-en-v1.5 & \texttt{d4aa690} \\ % numbers-ok: model name and revision hash
\gj{} members & PrismNLI-0.4B; bge-large-en-v1.5; bge-base-en-v1.5 & \texttt{02b9102}; \texttt{d4aa690}; \texttt{a5beb1e} \\ % numbers-ok: model names and revision hashes
Banking77 & dataset & \texttt{f541215} \\ % numbers-ok: revision hash
Bitext & dataset & \texttt{430d1a8} \\ % numbers-ok: revision hash
\bottomrule
\end{tabular}
\end{table}

\paragraph{Over-length requests.} Each model's release code defines what happens to a request that is too long,
and we applied that rule unchanged (\S\ref{sec:protocol}).
\begin{itemize}
  \item \pplx{} refuses a request above \dtLimitPplxRelease{} tokens; the refusal counts as wrong. Its serving
  engine, given the same limit, refused above \dtLimitPplxServed{} tokens; the one affected item was re-scored
  with the release rule (\S\ref{sec:find-serving}).
  \item \clef{} and \flash{} truncate the input state so that the request fits \dtLimitClefRelease{} tokens, the
  release code's default, and refuse only when the questions alone do not fit.
  \item \jev{} documents a limit for the state plus the longest question, and in our probe requests above it returned an
  error, not a truncated answer. No \jev{} request in this study came near it.
  \item \orb{} never truncates: its reader asserts that every prompt fits its context and stops otherwise.
\end{itemize}
No Banking77 or Bitext request reached any of these limits: the longest prompts we recorded had
\psMaxPromptTokens{} tokens on Banking77 and \btMaxPromptTokens{} on Bitext. The limits matter only in the thirteen-task study.
